% Complete unbounded closure, conditional only on the explicitly stated old mixed property.
% Original general conjecture remains open outside this reachable family.
\section{Full mixed closure below twice the first missing weight}
\begin{theorem}
Let $X$ be a normalized increasing arm with first missing
weight $\kappa\ge5$, quasi-centered deficit $e=e_X$, and
full mixed unimodality: $e q_Y$ is unimodal for every ordinary
increasing arm $Y$. Every legal append
\[
                   t=\kappa+g,\qquad0\le g\le\kappa-2
\]
preserves both quasi centering and full mixed unimodality.
The gap $g$ and the first missing weight are unbounded.
\end{theorem}

Write $c=\operatorname{wt}X$ and $D=c+1-\kappa=\deg e$.
The mass bound and $(\kappa-5)(\kappa-2)\ge0$ give
\[
                    D\ge2\kappa-4.                     \tag{1}
\]
In particular $t\le c$. The exact nonresonance theorem
therefore gives a quasi-centered new deficit
$f=[t]e+z^c$, of degree $c+g$ and layer endpoint sums
at most $c+g+1$.

For a nonrectangular opposite arm $Y$ write $b=\operatorname{wt}Y$,
$q=q_Y$, $p=e q$, $V=[t]p$. In the central degree orders
$b=c+s$ and at joining rank $c+j$, the actual increment is
\[
        A_j=p_{c+j}-p_{c+j-t}+a_j,\qquad
        a_0=1,\quad a_j=q_j-q_{j-1}\ (j\ge1).            \tag{2}
\]
Only $0\le j\le\lfloor(g+s+1)/2\rfloor\le g$ can
remain undecided by the new layer-center bound.

\paragraph{A strengthened resource classification, including the linear budget.}
In the central band $-g-1\le s\le g$, all increments of
the old $p$ from rank $c$ onward are nonpositive. Pairing
the full lowest deficit layer $[0,D]$ with ordinary layers
shows
\[
 p_{c+j-1}-p_{c+j}\ge q_{\kappa+j-2}-(F_Y)_{s-j}.
                                                               \tag{3}
\]
All other proper-layer products have nonpositive increments;
the full ordinary-layer window is nonincreasing there as well.

For $j\ge2$, the preceding section proves that the right side
pays every positive $a_j-a_{j-1}$, except possibly when an
ordinary long rectangle is already saturated. More precisely,
failure of this payment must belong to one of the following
two \emph{whole-profile} classes:
\begin{enumerate}
\item $Y=Z T$, $h=\operatorname{wt}Z<s$, and $T=b-h>c$.
\item $Y=1^\alpha U$, $H=\prod_{v\in U}[v]$ has
      degree $h\le s$, and $\alpha=b-\operatorname{wt}U>g$.
\end{enumerate}
To see the extra restrictions, in the first long-capacity case
let $H=q_Z$. At a saturated head $q_J=H(1)$ the diagonal is
$\sum_{i\le h-(s-j)}H_i-{\bf1}_{s=j}$.
If $h\ge s$ and $s>j$, it contains $H_j$.
If $s=j$, removing the empty-composition unit still leaves
at least $H_j$, since $H_0=1$ and $j\ge1$.
Either bound pays $H_j-H_{j-1}$. In the second case put
$P=\prod[v]$. Proper ordinary reflection gives
$(F_Y)_r\le\sum_{i\le r}P_i$ below the long initial interval.
If $\deg P\ge s+1$, symmetry makes the remaining tail
mass at least $P_j$, again paying the curvature.

The same classification holds for failure of the linear
budget with $j=1$, whose target is $a_1=q_1-1$.
Its bounded-capacity proof uses the main-box gap lemma at
$j=1$ and the identity $a_1=K_1-1$.
For a relevant first edge, $2\le g+s+1$ and the mass bound
imply $\deg K\ge\kappa$, so no small-degree exception occurs.
If $s-1<0$, $q_{\kappa-1}\ge K_1=q_1$ pays directly.
In a long-capacity case with $J<\operatorname{wt}Z$, the
CDF reflection lemma pays $H_1$ when $s\ge2$; at $s=1$
the initial cumulative sum includes $H_0+H_1$, cancelling
the one empty-composition unit. At a saturated head,
$\operatorname{wt}Z\ge s$ would again pay $H_1$.
The long-initial-ones argument likewise pays $P_1$ unless
$\deg P\le s$.

Thus either all the linear and adjacent resources pay, or
the fixed $Y$ belongs to one of the two displayed classes.
We handle the latter directly, without imposing a false
post-input-mode budget.

\paragraph{Direct joining concavity for the exceptional long rectangles.}
For a finite polynomial $P$, let
$C_P(n)=\sum_{i\le n}P_i$, with $C_P(n)=0$ for $n<0$.
Use a bar to reverse a polynomial in its own degree.

In class 1 put $H=q_Z$ of degree $h<s$. Throughout the
support $0\le i\le h$ of the correction, $c+i<b$ and
the exact window calculation gives
\[
 A_i=C_{\bar e\bar H}(h+g-i)
                -C_{eH}(h-s+i)+H_i.                   \tag{4}
\]
The extra final monomial in $q_Y=[T]H+z^b$ is absent at
these ranks. The first window is a tail cumulative sum
because $\deg(eH)=D+h<c$; the second is a head cumulative
sum because $h-s+i<t$.

In class 2 put $W=\operatorname{wt}U$ and
$H=\prod[v]$, of degree $h=W-|U|\le s$.
Define the ordinary padding kernel recursively by
\[
 K_\varnothing=1,\qquad K_{Zv}=1-z+z[v]K_Z.
\]
It is nonnegative: the product $[v]K_Z$ has constant
coefficient one, so the only negative source at rank one
cancels, and every higher coefficient is nonnegative.
Its degree is $W$ and its mass is $H(1)$.
The exact ordinary identity is
\[
                       (1-z)q_Y=H-z^{\alpha+1}\bar K_U.
\]
Now $a_i=H_i$ through the joining band, and
\[
 A_i=C_{\bar e\bar H}(h+g-i)
                  -C_{e\bar K_U}(W+i-s-1)+H_i
                                    \qquad(0\le i\le h).       \tag{5}
\]
Again the first term is a tail cumulative sum. The second
index is smaller than $t$, since $W\le2h$, $h\le s\le g$
and $g\le\kappa-2$.

In either formula, for $1\le i\le h$ the first cumulative
term loses the coefficient
$(\bar e\bar H)_{h+g-i+1}$, while the second loses a
nonnegative coefficient. By (1),
\[
                  h\le h+g-i+1\le2g\le D.
\]
Every coefficient of $\bar e$ in its support is a positive
integer. Hence this first loss is at least $H(1)$: all
coefficients of $\bar H$ are sampled with multiplicity at
least one. This pays $H_i-H_{i-1}$, proving
\[
                         A_i\le A_{i-1}\qquad(1\le i\le h).
                                                               \tag{6}
\]
For $i>h$ the correction increment is zero, so $A_i$ is
an increment of the unimodal old window $V$. If any
earlier $A_i$ is negative, then the corresponding increment
of $V$ is negative as well, because $H_i\ge0$.
Unimodality of $V$ prevents any later positive increment
once the correction has ended. Thus (6) proves the entire
joining no-valley statement in both exceptional classes.

The unbounded coefficient-mass domination, cumulative-slope
implication, support arithmetic and zero-correction splice
are kernel checked in \texttt{LongRectangleRepair.lean}.
The exact ordinary-polynomial identifications (4)--(5)
are the written calculations above, independently checked
with arbitrary positive integer deficit profiles in
\texttt{audit\_long\_rectangle\_repair.json}.

\paragraph{Closing the full mixed induction.}
Induct on the number of non-one parts of $Y$; rectangular
opposite arms are uniform-window bases.
If $b\ge c+g+1$, deleting the last part of $Y$ gives
the old new-deficit mixed window below rank $b$.
The quasi-center bound makes every proper ordinary-layer
product nonincreasing at $b$. The full ordinary layer is
constant there because it covers the complete support of $f$.
Thus the induction prefix joins a nonincreasing tail.
If $b\le c-g-2$, the center bound gives the same join at
rank $c$, before which the shifted correction is absent.

In the remaining band $c-g-1\le b\le c+g$, the established
first-edge dichotomy says $p_c<p_{c-1}$ or $p_c=\max p$.
Integer unit repair handles $A_0$ and the earlier window
prefix. If all the resources in (3) pay, then
\[
 p_c-p_{c+j}\ge a_j,\qquad
 p_{c+j-1}-p_{c+j}\ge a_j-a_{j-1}.
\]
Indeed the first linear budget starts the telescoping sum;
nonpositive curvatures need no payment because $p$ is
nonincreasing there. The general multi-edge repair therefore
propagates every positive joining increment backwards.
If any resource fails, (4)--(6) give the same no-valley
conclusion directly. Beyond the last joining edge, the
new center bound makes all actual increments nonpositive.
This proves full mixed unimodality of $f$ and closes the
induction on $Y$.

The argument combines the established unit/multi-edge repair
theorems with the new exact resource and long-rectangle
lemmas. It does not infer coefficient unimodality from
cumulative log-concavity.

\begin{corollary}[Arbitrarily many doubling phases]
All arms
\[
 1^\alpha4^{m_4}8^{m_8}16^{m_{16}}\cdots(2^L)^{m_{2^L}}
\]
have full mixed unimodality against arbitrary opposite
increasing arms if $\alpha\ge1$, $m_4\ge3$, every nonfinal
phase $u\ge8$ has $m_u\ge u/2$, and the final phase has
at least one copy. Their strict-peak cores and both
relative polynomials are unimodal by the established
short-prefix reflection join.
\end{corollary}
\begin{proof}
The three-four seed has first missing weight five and is
already quasi centered with full mixed unimodality.
After a saturated phase $u$, the next capacity $2u$ is
exactly $2\kappa-2$, where $\kappa=u+1$.
Its next $u-1$ repeats raise first missing weight to $2u+1$;
all those appends fall in the same new range. Further
repeats use the established stationary closure.
Iterate the theorem and apply the existing core join.
\end{proof}

This completes a previously open unbounded range of append
gaps. It leaves the general unamplified conjecture open:
large jumps $t\ge2\kappa-1$, small old first missing weights,
and general non-quasi-centered prefixes are not covered.
